\documentclass[12pt]{article}
\title{A Sample \LaTeX\ Document}
\author{Math 300}
\usepackage{graphicx}
\usepackage{mathdots} % para el comando \iddots
\usepackage{amssymb, amsmath, amsbsy} % simbolitos
\usepackage{xcolor}

\begin{document}
%\maketitle

pag18

\hfill

\textbf{Chapter 1
Examples and definitions of dynamical
phenomena
}

\hfill

A dynamical system can be any mechanism that evolves deterministically in time.
Simple examples can be found in mechanics, one may think of the pendulum or the
solar system. Similar systems occur in chemistry and meteorology.We should note
that in biological and economic systems it is less clear when we are dealing with
determinism.1 We are interested in evolutions, that is, functions that describe the
state of a system as a function of time and that satisfy the equations of motion of the
system. Mathematical definitions of these concepts follow later.

\hfill

The simplest type of evolution is stationary, where the state is constant in time.
Next we also know periodic evolutions. Here, after a fixed period, the system always
returns to the same state and we get a precise repetition of what happened in the
previous period. Stationary and periodic evolutions are very regular and predictable.
Here prediction simply is based on matching with past observations. Apart from
these types of evolutions, in quite simple systems one meets evolutions that are not
so regular and predictable. In cases where the unpredictability can be established in
a definite way, we speak of chaotic behaviour.

\hfill

In the first part of this chapter, using systems such as the pendulum with or without
damping or external forcing, we give an impression of the types of evolution
that may occur. We show that usually a given system has various types of evolution.
Which types are typical or prevalent strongly depends on the kind of system
at hand; for example, in mechanical systems it is important whether we consider
dissipation of energy, for instance by friction or damping. During this exposition
the mathematical framework in which to describe the various phenomena becomes
clearer. This determines the language and the way of thinking of the discipline of
Dynamical Systems.

\hfill

In the second part of this chapter we are more explicit, giving a formal definition
of dynamical system. Then concepts such as state, time, and the like are also
discussed. After this, returning to the previous examples, we illustrate these concepts.
In the final part of the chapter we treat a number of miscellaneous examples
that regularly return in later parts of the book. Here we particularly pay attention to
examples that, historically speaking, have given direction to the development of the
discipline. These examples therefore act as running gags throughout the book.

\hfill

\textbf{1.1 The pendulum as a dynamical system}

\hfill

As a preparation to the formal definitions we go into later in this chapter, we first
treat the pendulum in a number of variations: with or without damping and external
forcing.We present graphs of numerically computed evolutions. From these graphs
it is clear that qualitative differences exist between the various evolutions and that
the type of evolution that is typical or prevalent strongly depends on the context
at hand. This already holds for the restricted class of mechanical systems considered
here.

\hfill

\textbf{1.1.1 The free pendulum}

\hfill

The planar mathematical pendulum consists of a rod, suspended at a fixed point in
a vertical plane in which the pendulum can move. All mass is thought of as being
concentrated in a point mass at the end of the rod (see Figure 1.1), and the rod itself
is massless. Also the rod is assumed stiff. The pendulum has mass m and length `.
We moreover assume the suspension to be such that the pendulum not only can
oscillate, but also can go ‘over the top’. In the case without external forcing, we
speak of the free pendulum.

\hfill

\textbf{1.1.1.1 The free undamped pendulum}

\hfill


In the case without damping and forcing the pendulum is only subject to gravity,
with acceleration g: The gravitational force is pointing vertically downward with
strength mg and has a component $-mg \sin \varphi$ along the circle described by the
point mass; see Figure 1.1. Here $\varphi$ is the angle between the rod and the downward
vertical, often called ‘deflection’, expressed in radians. The distance of the point
mass from the ‘rest position’ ($\varphi = 0$), measured along the circle of all its possible
positions, then is $l \varphi$. The relationship between force and motion is determined by
Newton’s2 famous law F D ma; where F denotes the force, m the mass, and a the
acceleration.

\hfill

By the stiffness of the rod, no motion takes place in the radial direction and we
therefore just apply Newton’s law in the $\varphi$-direction. The component of the force in
the $\varphi$-direction is given by $-mg \sin \varphi$, where the minus sign accounts for the fact
that the force is driving the pendulum back to $\varphi$. For the acceleration $a$ we have

\[
a = \frac{d^2 (l \varphi)}{dt^2} = l \frac{d^2 \varphi}{dt^2}
\]

where $(d^2 \varphi / dt^2)(t) = \varphi''(t)$ is the second derivative of the function $t \to \varphi(t)$. Substituted into Newton’s law this gives

\[
ml \varphi '' = -mg \sin \phi
\]

or, equivalently,

\[
\varphi '' = -\frac{g}{l} \sin \varphi
\]


where obviously $m, l > 0$. So we derived the equation of motion of the pendulum,
which in this case is an ordinary differential equation. Thismeans that the evolutions
are given by the functions $t \mapsto \varphi(t)$ that satisfy the equation of motion (1.1). In the
sequel we abbreviate $w = \sqrt{g/l}$.

\hfill

Observe that in the equation of motion (1.1) the mass m no longer occurs. This
means that the mass has no influence on the possible evolutions of this system.
According to tradition, as an experimental fact this was already known to Galileo,3
a predecessor of Newton concerning the development of classical mechanics.

\hfill

Remark (Digression on Ordinary Differential Equations I). In the above example,
the equation of motion is an ordinary differential equation. Concerning the solutions
of such equations we make the following digression.


\begin{enumerate}
\item From the theory of ordinary differential equations (e.g., see [16, 69, 118, 144])
it is known that such an ordinary differential equation, given initial conditions
or an initial state, has a uniquely determined solution.4 Because the differential
equation has order two, the initial state at $t=0$ is determined by the two
data $\varphi(0)$ and $\varphi'(0)$, thus both the position and the velocity at the time $t=0$.
The theory says that, given such an initial state, there exists exactly one solution
$t \mapsto \varphi(t)$, mathematically speaking a function, also called ‘motion’. This
means that position and velocity at the instant t= 0 determine the motion for
all future time.5 In the discussion later in this chapter, on the state space of a
dynamical system, we show that in the present pendulum case the states are
given by pairs $(\varphi, \varphi')$. This implies that the evolution, strictly speaking, is a map
$t \mapsto (\varphi(t), \varphi'(t))$, where $t \mapsto \varphi(t)$ satisfies the equation of motion (1.1). The
plane with coordinates $(\varphi, \varphi')$ often is called the phase plane. The fact that the
initial state of the pendulum is determined for all future time, means that the
pendulum system is deterministic.6

\item Between existence and explicit construction of solutions of a differential equation
there is a wide gap. Indeed, only for quite simple systems,7 such as the harmonic
oscillator with equation of motion

\[
\varphi'' = -w^2 \varphi,
\]

can we explicitly compute the solutions. In this example all solutions are of the
form $\varphi(t) =  A \cos(w t + B)$, where A and B can be expressed in terms of the
initial state $(\varphi(0), \varphi'(0))$. So the solution is periodic with period $2\pi/w$ and with
amplitude A; whereas the number B; which is only determined up to an integer
multiple of $2\pi$, is the phase at the time t = 0: Because near $\varphi = 0$ we have the
linear approximation

\[
\sin \varphi \approx \varphi
\]

in the sense that $\sin \varphi = \varphi+ o(|\varpi|)$, we may consider the harmonic oscillator
as a linear approximation of the pendulum. Indeed, the pendulum itself also has
oscillating motions of the form $varphi(t) = F(t)$, where F(t) = F(t+P) for
a real number P; the period of the motion. The period P varies for different
solutions, increasing with the amplitude. It should be noted that the function F
occurring here, is not explicitly expressible in elementary functions. For more
information see [153].

\item General methods exist for obtaining numerical approximations of solutions of
ordinary differential equations given initial conditions. Such solutions can only
be computed for a restricted time interval. Moreover, due to accumulation of
errors, such numerical approximations generally will not meet with reasonable
criteria of reliability when the time interval is growing (too) long.

\hfill

We now return to the description of motions of the pendulum without forcing or
damping, that is, the free undamped pendulum. In Figure 1.2 we present a number
of characteristic motions, or evolutions, as these occur for the pendulum.

\hfill

In diagrams (a) and (b) we plotted $\varphi$ as a function of t: Diagram (a) shows an
oscillatory motion, that is, where there is no going ‘over the top’ (i.e., over the
point of suspension). Diagram (b) shows a motion where this does occur. In the
latter case the graph does not directly look periodic, but when realising that $\varphi$ is
an angle variable, which means that only its values modulo integer multiples of $2\phi$ 
are of interest for the position, we conclude that this motion is also periodic. In
diagram (c) we represent motions in a totally different way. Indeed, for a number of
possible motions we depict the integral curve in the $(\varphi, \varphi')$-plane, that is, the phase
plane. The curves are of the parametrised form $t \mapsto (\varphi(t), \varphi'(t)),$ where t varies
over $\mathbb{R}$: We show six periodic motions, just as in diagram (a), noting that in this
representation each motion corresponds to a closed curve. Also periodicmotions are
shown as in diagram (b), for the case where $varphi$ continuously increases or decreases.
Here the periodic character of the motion does not show so clearly. This becomes
more clear if we identify points on the $\varphi$-axis which differ by an integer multiple
of $2\phi$.

\hfill

Apart from these periodic motions, a few other motions have also been represented.
The points $(2\pik,0)$, with integer k; correspond to the stationary motion
(which rather amounts to rest and not to motion) where the pendulum hangs in
its downward equilibrium. The points .2 .k C 1
2 /; 0/ k 2 Z correspond to the
stationary motion where the pendulum stands upside down. We see that a stationary
motion is depicted as a point; such a point should be interpreted as a ‘constant’
curve. Finally motions are represented that connect successive points of the form
.2 .k C 1
2 /; 0/ k 2 Z: In these motions for t ! 1 the pendulum tends to
its upside-down position, and for t !  1 the pendulum also came from this
upside-down position. The latter two motions at first sight may seem somewhat
exceptional, which they are in the sense that the set of initial states leading to such
evolutions have area zero in the plane. However, when considering diagram (c) we
see that this kind of motion should occur at the boundary which separates the two
kinds of more common periodicmotion. Therefore the corresponding curve is called
a separatrix. This is something that will happenmore often: the exceptionalmotions
are important for understanding the ‘picture of all possible motions.’
Remark. The natural state space of the pendulum really is a cylinder. In order to see
that the latter curves are periodic, we have to make the identification '   ' C 2 ;
which turns the .'; '0/-plane into a cylinder. This also has its consequences for
the representation of the periodic motions according to Figure 1.2(a); indeed in the
diagram corresponding to (c) we then would only witness two of such motions. Also
we see two periodic motions where the pendulum goes ‘over the top’.
For the mathematical understanding of these motions, in particular of the representations
of Figure 1.2(c), it is important to realise that our system is undamped.
Mechanics then teaches us that conservation of energy holds. The energy in this case
is given by
H.'; '0/ D m`2  1
2
.'0/2   !2 cos ' : (1.2)
Notice that H has the format ‘kinetic plus potential energy.’ Conservation of energy
means that for any solution '.t/ of the equation of motion (1.1), the function
H.'.t/; '0.t // is constant in t: This fact also follows from more direct considerations;
compare Exercise 1.3. This means that the solutions as indicated in diagram
(c) are completely determined by level curves of the function H: Indeed, the level
curves with H-values between  m`2!2 and Cm`2!2 are closed curves corresponding
to oscillations of the pendulum. And each level curve withH > m`2 !2
consists of two components where the pendulum goes ‘over the top’: in one component
the rotation is clockwise and in the other one counterclockwise. The level
H D m`2!2 is a curve with double points corresponding to the exceptionalmotions
just described: the upside-down position and the motions that for t !˙1tend to
this position.
We briefly return to a remark made in the above digression, saying that the
explicit expression of solutions of the equation of motion is not possible in terms of
elementary functions. Indeed, considering the level curve with equation H.'; '0/ D E; we solve to
'0 D ˙
s
2   E
m`2 C !2cos' :
\end{enumerate}



\begin{thebibliography}{X}

\bibitem{Broer} Henk Broer, Floris Takens, {\bf Dynamical Systems
and Chaos,} Springer, 2009.


\end{thebibliography}
\end{document}

